\documentclass{article}
\usepackage{amsmath, amssymb, amsfonts}
\usepackage{geometry}
\geometry{margin=1in}

\begin{document}

\noindent\textbf{(i)} given basis of topology has $k$ elements.

$\Rightarrow$ by def\textsuperscript{n} of basis. all open set in that shall could be written as union of some or all of these basis's.

so. given $k$. basis.

\quad we have $2^k - 1$ (non empty union of basis's) giving $2^k - 1$ open set that are not empty, including whole set $(X)$ too.

$\therefore$ topology can at most have all possible open set including \& of empty set.

\quad giving us $2^k - 1 + 1$ $(\phi)$ open set

\quad i.e. total of $2^k$ open sets.

\vspace{1em}

\noindent\textbf{(ii)} On given set, $X$, sub-basis has $n$ elements there by def\textsuperscript{n} of sub-basis, basis elements can be formed of intersections (finite in this case) of some or all of the elements of sub-bases.

\quad giving us at most $(2^n)$ basis.

\& by Part 1.

\quad our set will have at most

\quad $2^{(2^n)}$ open sets.

\end{document}